\pdfinfoomitdate=1
\pdftrailerid{}
\documentclass[10pt]{article}
\usepackage[margin=0.82in]{geometry}
\usepackage[T1]{fontenc}
\usepackage{lmodern,amsmath,amssymb,amsthm,booktabs,microtype}
\usepackage[hidelinks]{hyperref}
\hypersetup{pdftitle={A past-future scalar comparison for a timelike-shell packet},pdfauthor={Ricardo Maldonado}}
\newtheorem{theorem}{Theorem}
\newtheorem{corollary}[theorem]{Corollary}
\newcommand{\Hh}{\mathcal H}
\newcommand{\dd}{\Delta}
\newcommand{\norm}[1]{\lVert#1\rVert}
\title{A past--future scalar comparison for a timelike-shell packet}
\author{Ricardo Maldonado\\\small Independent Researcher\\\small Henderson, Nevada, United States}
\date{4 November 2026 (America/Los\_Angeles)}
\begin{document}
\maketitle
\begin{abstract}
For one registered quadrupole packet in a conditional linear Einstein--scalar
shell construction, the complete dynamical boundary problem is even in time,
although the initial velocity is nonzero. Comparing the developments with
opposite initial velocities removes the initial displacement from a higher
graph-energy estimate. This gives the uniform absolute bound
\( |S(-t)-S(t)|<264000000001/13738809600<20 \) for
\(0\le t\le3\ell\), \(\ell=10^{-6}\), where \(S\) is the first physical
scalar amplitude on the shell in fixed model units. After an explicit
state, clock and normalization join to retained future waveform enclosures,
two lobe signs first give a coarse past zero. Two selected whole-cell
enclosures then localize at least one amplitude zero to
\((-6145\ell/4096,-6143\ell/4096)\), a gap of width \(\ell/2048\).
The result retains the earlier branch, reconstruction,
regularity, uniqueness and waveform-certificate premises. It supplies neither
a nonlinear amplitude interval nor evidence for a cosmological origin.
\end{abstract}

\section{Registered problem and conditional result}

Fix \(\ell=10^{-6}\), \(x=t/\ell\), \(r=z/\ell\), and the flat unit bump
\[
 Q(r)=\begin{cases}\exp\!\left(4-\dfrac1{(r-1)(2-r)}\right),&1<r<2,\\0,&\text{otherwise}.\end{cases}
 \qquad g(r)=\ell e^{-\ell r}Q(r).
\]
The registered master data are \(U(0,r)=Q(r)\), \(U_x(0,r)=g(r)\).
The physical first scalar coefficient, including the normal embedding in
the retained straight-shell reconstruction, is
\begin{equation}\label{eq:observation}
 \Phi_1(\tau,\Omega)=S(\tau/\rho_b)Y(\Omega),\qquad
 Y=n_1^2-n_2^2,\quad \tau=\rho_b t.
\end{equation}
Here \(\sup|Y|=1\) and its angular mean is zero. The background shell scalar
is constant, so first clock/label pullbacks add no amplitude term.

We retain the analytic branch, admitted shell enclosures, smooth linear
development, full dynamical junction law and physical reconstruction of
the past-contact construction~\cite{contact}. We also identify explicitly
when the older actual future waveform enclosure~\cite{future} is used.
Selected-source hashing and finite arithmetic do not reprove these premises.

\begin{theorem}\label{thm:main}
Under those retained premises, the actual registered first scalar amplitude
satisfies
\begin{equation}\label{eq:main}
 \sup_{0\le t\le3\ell}|S(-t)-S(t)|<B,
 \qquad B=\frac{264000000001}{13738809600}<20.
\end{equation}
The compared proper times are \(-\rho_b t\) and \(+\rho_b t\) at the same
anchor. The bound concerns a first linear coefficient in model units.
\end{theorem}

This is an absolute difference estimate: it does not bound either individual
amplitude by 20, prescribe a relative error at a zero, or assert equality of
past and future profiles. The construction uses established time-reflection,
positive-form and graph-energy ideas~\cite{zahn,marta}; its contribution is
the explicit estimate and normalized certificate transfer for this packet.

\section{A common operator and a zero-displacement difference}

Use the same auxiliary static half-line realization for both developments:
\begin{align}
 \Hh&=L^2(0,\infty)\oplus\mathbb R,
 &\norm{(u,q)}_{\Hh}^2&=\int_0^\infty u^2\,dr+\Lambda q^2,\label{eq:space}\\
 A_0(u,q)&=\left(-u_{rr}+V_0u,\frac{-u_r(0)+b_0q}{\Lambda}\right),
 &\Lambda&=\frac{\mu_b}{\ell}>1000,\label{eq:operator}\\
 V_0&=\ell^2\rho^2Q_{\rm bg},
 &b_0&=\ell\left(\gamma+\frac52\mu_b\right),\quad\gamma=\frac{C_b}{2}.
 \nonumber
\end{align}
The form domain links \(q=u(0)\); the operator domain has no extra static
Robin condition. Both endpoint kinetic and spring terms are retained,
as in the standard coupled energy bookkeeping~\cite{dahmani}.
The prior positive-form construction gives \(A_0\ge5500\ell^2 I\).
The complete equation and dynamic boundary law take the common form
\begin{equation}\label{eq:wave}
 Z_{xx}+(A_0+\omega(x)I)Z=0,\qquad
 \omega(x)=-\frac{35}{4}\ell^2\tanh^2(\ell x),\qquad
 b(x)=b_0+\Lambda\omega(x).
\end{equation}
In particular, the boundary spring has the same time dependence as the bulk
operator; omitting that endpoint term would invalidate the argument.

The static extension agrees with the actual coefficients through \(r=6\),
within the validated background collar \(r\le8\).
The original fields have support within \([0,2+s]\), and before \(s=1\)
within \([1-s,2+s]\), for \(s=|x|\le3\). They therefore stay in \(r\le5\),
where the extension has no effect. Their flat initial data satisfy every
power-domain compatibility requirement in the retained smooth construction.

Let \(Z_+\) have data \((Q,+g)\) and \(Z_-\) data \((Q,-g)\), including
the linked zero initial endpoint components. Evenness of the \emph{complete}
equation and boundary law, followed by uniqueness, gives
\[
 Z_{\rm current}(s)=Z_+(s),\qquad Z_{\rm current}(-s)=Z_-(s).
\]
Thus \(\dd Z=Z_--Z_+\) has \(\dd Z(0)=0\),
\(\dd Z_s(0)=(-2g,0)\). The factor two in the velocity is essential.
For positive \(s\), \(\omega_s\le0\), so the complete estimating energy
\begin{equation}\label{eq:energy}
 E[X]=\frac12\left(\norm{X_s}_{\Hh}^2+
       \norm{A_0^{1/2}X}_{\Hh}^2+\omega\norm X_{\Hh}^2\right),
 \qquad E'[X]=\frac{\omega_s}{2}\norm X_{\Hh}^2\le0
\end{equation}
is nonincreasing. It is a mathematical norm estimate, with no physical
injected-energy or Noether-energy interpretation supplied here.
Since \(\norm Q_2\le1\),
\(E[\dd Z](s)\le E[\dd Z](0)\le2\ell^2\).
Its local endpoint trace \(\dd q\) vanishes before contact. From the endpoint
mass in~\eqref{eq:energy}, \(500|\dd q_s|^2\le E[\dd Z](0)\), whence
\begin{equation}\label{eq:qbound}
 |\dd q(s)|^2\le\frac{(s-1)_+^2}{500}E[\dd Z](0)
 \le\frac{2\ell^2}{125}<\ell^2,\qquad0\le s\le3.
\end{equation}

\section{Graph energy and the physical scalar bound}

Apply the abstract form identity to \(A_0^{1/2}\dd Z\). This vector is
generally nonlocal; no compact-support assertion is made for it.
Scalar \(\omega I\) commutes with \(A_0\), and coercivity gives
\begin{align}
 2E[A_0^{1/2}\dd Z](s)&\ge(1-\epsilon)\norm{A_0\dd Z(s)}_{\Hh}^2,
 &\epsilon&=\frac{63\ell^2}{4400}<1,\label{eq:coercive}\\
 2E[A_0^{1/2}\dd Z](0)&=4a_0(g,g)
 \le4\ell^2\bigl[(16+\ell)^2+103000\ell^2\bigr].\label{eq:graphinitial}
\end{align}
To see the constants, \(|\omega|\le(315/4)\ell^4\) for \(s\le3\),
and \(V_0<103000\ell^2\) on the initial support. Writing
\(y=[(r-1)(2-r)]^{-1}\ge4\) gives
\(Q_r=(3-2r)y^2e^{4-y}\), so \(|Q_r|\le16\) analytically and
\(\norm{g_r}_2\le\ell(16+\ell)\).
The zero displacement removes the initial \(Q_{rr}\) contribution.
The spectral identity in~\eqref{eq:graphinitial} includes every endpoint
component of the half-power vector. Exact arithmetic yields
\begin{equation}\label{eq:graphbound}
 \norm{A_0\dd Z(s)}_{\Hh}^2
 \le\frac{4\ell^2[(16+\ell)^2+103000\ell^2]}{1-\epsilon}
 <(33\ell)^2.
\end{equation}

The retained physical reconstruction and boundary cancellation give
\begin{equation}\label{eq:physical}
 S_{\dd}(\ell s)=\frac{3\mu_b\cosh(\ell s)^{-3/2}}{\rho_b^3p_b}
 \left[-\frac{(A_0\dd Z)_b}{\ell^2}+\frac52\dd q\right].
\end{equation}
Indeed \(U_r=b_0q-\Lambda(A_0Z)_b\) implies
\(U_r/\ell-\gamma q=\mu_b[-(A_0Z)_b/\ell^2+(5/2)q]\).
The temporal term cancels. The normal derivative is spatial and the
conjugating \(\cosh\) factor is even, so there is no additional reflection
minus sign: \(S_\dd(t)=S(-t)-S(t)\).

The endpoint norm implies
\(\mu_b|(A_0\dd Z)_b|\le\sqrt{\mu_b\ell}\norm{A_0\dd Z}_{\Hh}\).
Use the admitted bounds
\(\mu_b<1/400\), \(\sqrt{\mu_b\ell}<1/20000\),
\(\rho_b\ge129\), \(|p_b|\ge3/25000\), together with
\eqref{eq:qbound} and~\eqref{eq:graphbound}. Then
\begin{equation}\label{eq:ceiling}
 |S_\dd|<\frac{25000}{129^3}
 \left(\frac{33}{20000\ell}+\frac{\ell}{160}\right)
 =\frac{264000000001}{13738809600}=B<20.
\end{equation}
This proves Theorem~\ref{thm:main}. It also bounds the angular first-field
difference because \(|Y|\le1\), without supplying an SI calibration.

\section{A retained future certificate and a coarse past zero}

The selected wider inertial waveform proof~\cite{future} uses this same
\(\ell=10^{-6}\), branch, flat bump, nonzero velocity \(e^{-z}Q(z/\ell)\),
quadrupole, conjugation and full dynamic shell law. Its actual enclosure
includes the proved error from the free comparator. The report exports
\begin{equation}\label{eq:export}
 H(x)=\ell\,[w_z(\ell x,0)-\gamma w(\ell x,0)],\qquad
 \ell S(\ell x)=\frac{3\cosh(\ell x)^{-3/2}}{\rho_b^3p_b}H(x).
\end{equation}
Thus the stored \texttt{ell\_S} interval must be divided by \(\ell\)
\emph{before} adding the model-unit difference ceiling \(B\).
The same reconstructed state and clock are identified by retained linear
uniqueness; neither a comparator alone nor the separate \(\ell=10^{-10}\)
packet supplies this join.

If \([L,U]\) encloses the actual future value \(S(t)\), then
\([L-B,U+B]\) encloses the matched past value \(S(-t)\).
Table~\ref{tab:transfer} shows outward decimal roundings of the exact
rational endpoints. The proof uses the rational intervals, not those
readable decimal summaries.
\begin{table}[ht]
\centering
\small
\begin{tabular}{@{}lll@{}}
\toprule
Matched \(t>0\)&Retained future \(S(t)\)&Transported past \(S(-t)\)\\
\midrule
\(5\ell/4\)&\([-41241.437,-37396.209]\)&\([-41260.653,-37376.993]\)\\
\(7\ell/4\)&\([37398.637,41244.115]\)&\([37379.422,41263.331]\)\\
\bottomrule
\end{tabular}
\caption{Two actual coefficient enclosures in model units, conditional on
the retained future enclosure theorem and Theorem~\ref{thm:main}.}
\label{tab:transfer}
\end{table}

\begin{corollary}
With the retained future certificate, \(S(-5\ell/4)<-37376\) and
\(S(-7\ell/4)>37379\). The continuous amplitude therefore has at least
one zero in \((-7\ell/4,-5\ell/4)\).
\end{corollary}

\paragraph{A selected-cell refinement.}
The older certificate also supplies continuous enclosures on whole closed
grid cells, rather than only sampled values. Authenticating its cell CSV
and retaining the earlier full-cell enclosure theorem gives the two
uniform transports
\begin{equation}\label{eq:cells}
 \begin{array}{c|c}
 s\in[3071/2048,\,6143/4096]&S(-\ell s)<-53\\[2pt]
 s\in[6145/4096,\,3073/2048]&S(-\ell s)>63.
 \end{array}
\end{equation}
The CSV stores integer \texttt{ell\_S} endpoints with value denominator
\(D_v=2^{112}\). Decode them by \(D_v\), then divide by \(\ell\) to obtain
physical-\(S\) endpoints and widen by \(B\). No old waveform
campaign or numerical PDE is rerun. Applying continuity to the two inner
cell endpoints proves the sharper conditional statement
\begin{equation}\label{eq:refinedzero}
 \exists\,t_0\in
 \left(-\frac{6145\ell}{4096},-\frac{6143\ell}{4096}\right):
 \ S(t_0)=0,\qquad
 \text{gap width}=\frac{\ell}{2048}.
\end{equation}
This gap is exactly 1024 times narrower than the preceding coarse past
interval of width \(\ell/2\). This comparison concerns those two past
intervals only; it does not compare precision with older fine future
brackets or assert optimal localization.

Neither zero conclusion gives uniqueness or a total zero count. These are
signs of the coefficient \(S\), rather than a globally signed angular field.
The four older narrow central/late point witnesses and negative future tail
have margins smaller than \(B\); their widened intervals straddle zero.
Those particular older point-witness brackets and the tail sign therefore
do not transfer by this ceiling. The selected-cell argument establishes a
different past interval. A quadratic response integral needs a separate norm and
cross-term argument, rather than pointwise interval widening.

\paragraph{Scope and reproducibility.}
Selected local input copies, a continuous primary proof, an independent
odd-part derivation, a separately checked whole-cell certificate join,
exact normal/optimized arithmetic reports and the
manuscript display audit accompany this note in the
\href{https://github.com/maldonado-research/HDblast/tree/main/research/HDBLAST_PAST_FUTURE_SCALAR_TRANSFER_20261104}{research capsule}.
The retained older waveform theorem is accepted as a conditional premise;
its full campaign is not rerun. The new calculations use rational arithmetic
and source checks, with no numerical PDE or field/source/bump oracle.
Internal review is distinct from external peer review and formal
proof-assistant certification.

Earlier future-contact work~\cite{second} already constructs a conditional complete
finite-angular second Taylor coefficient under its correction, gauge,
corner and support hypotheses. This note concerns the first coefficient.
Quantitative control of the chosen past second-order observable and
particular baseline, a nonlinear defect-to-error map and effective parameter
interval, physical incoming preparation, radiation and cosmological-origin
evidence remain separate tasks.

\begin{thebibliography}{9}\small
\bibitem{contact} R.~Maldonado, \emph{A pointwise scalar-trace bound through
past shell contact}, working paper and source capsule (4 November 2026),
\href{https://github.com/maldonado-research/HDblast/tree/b2a2fce15ca1870fff3e4bc312a65480c9e839cf/research/HDBLAST_PAST_SHELL_CONTACT_20261104}{HDBLAST past-shell-contact checkpoint}.
\bibitem{future} R.~Maldonado, \emph{Inertial shell waveform}, primary proof,
actual enclosure report and selected source inputs (15 October 2026),
\href{https://github.com/maldonado-research/HDblast/tree/415ab642d0a13720c03594c6968f9aede67d4c34/research/HDBLAST_INERTIAL_SHELL_WAVEFORM_20261015}{HDBLAST inertial-waveform checkpoint}.
\bibitem{zahn} J.~Zahn, \emph{Generalized Wentzell boundary conditions and
quantum field theory}, Annales Henri Poincar\'e \textbf{19}, 163 (2018),
\href{https://doi.org/10.1007/s00023-017-0629-3}{doi:10.1007/s00023-017-0629-3};
\href{https://arxiv.org/abs/1512.05512}{arXiv:1512.05512}.
\bibitem{marta} A.~Marta, \emph{The nonlinear wave equation with nonlinear
Wentzell boundary conditions on time-dependent compact Riemannian manifolds},
\href{https://arxiv.org/abs/2404.05855}{arXiv:2404.05855} (2024),
especially Remark~4.3 on reversal of time and Remarks~4.4--4.5 on the
coupled energy and finite speed.
\bibitem{dahmani} A.~Dahmani, Y.~Chitour, H.-M.~Nguyen and C.~Roman,
\emph{One dimensional wave equation with in-domain localized damping and
Wentzell boundary conditions},
\href{https://arxiv.org/abs/2503.22439}{arXiv:2503.22439} (2025),
Section~4, Eqs.~(4.2)--(4.3). Its damping-dependent decay estimate is not
imported here.
\bibitem{second} R.~Maldonado, \emph{A complete finite-angular second-order
coefficient through shell contact} (14 October 2026),
\href{https://github.com/maldonado-research/HDblast/blob/f4001e5d94e6388bfff4f82dee63c174f2dfd5eb/research/HDBLAST_FULL_SECOND_ORDER_CONTACT_20261014/FULL_FINITE_ANGULAR_SECOND_ORDER_CONTACT.md}{HDBLAST full-second-order-contact primary proof}, Sections~1 and~8.
\end{thebibliography}
\end{document}
